% problem_id: S001-lemma-algebraization-principal
% source_id: S001 (Stacks Project)
% source_file: algebraization.tex
% statement_locator: algebraization.tex lines 6236--6255
% proof_locator: algebraization.tex lines 6257--6365
% env: lemma
% pairing: tight (verified)
% status: complete (statement + author proof verbatim + reason + locators)
%
\begin{lemma}
\label{lemma-algebraization-principal}
In Situation \ref{situation-algebraize} let $(\mathcal{F}_n)$ be an object
of $\textit{Coh}(U, I\mathcal{O}_U)$. Assume
\begin{enumerate}
\item $A$ is local and $\mathfrak a = \mathfrak m$ is the maximal ideal,
\item $A$ has a dualizing complex,
\item $I = (f)$ is a principal ideal for a nonzerodivisor $f \in \mathfrak m$,
\item $\mathcal{F}_n$ is a finite locally free
$\mathcal{O}_U/f^n\mathcal{O}_U$-module,
\item if $\mathfrak p \in V(f) \setminus \{\mathfrak m\}$, then
$\text{depth}((A/f)_\mathfrak p) + \dim(A/\mathfrak p) > 1$, and
\item if $\mathfrak p \not \in V(f)$ and
$V(\mathfrak p) \cap V(f) \not = \{\mathfrak m\}$, then
$\text{depth}(A_\mathfrak p) + \dim(A/\mathfrak p) > 3$.
\end{enumerate}
Then $(\mathcal{F}_n)$ extends canonically to $X$. In particular, if $A$
is complete, then $(\mathcal{F}_n)$ is the completion of a coherent
$\mathcal{O}_U$-module.
\end{lemma}

\begin{proof}
We will prove this by verifying hypotheses (a), (b), and (c) of
Lemma \ref{lemma-when-done}.

\medskip\noindent
Since $\mathcal{F}_n$ is locally free over $\mathcal{O}_U/f^n\mathcal{O}_U$
we see that we have short exact sequences
$0 \to \mathcal{F}_n \to \mathcal{F}_{n + 1} \to \mathcal{F}_1 \to 0$
for all $n$. Thus condition (b) holds by
Cohomology, Lemma \ref{cohomology-lemma-topology-I-adic-f}.

\medskip\noindent
By induction on $n$ and the short exact sequences
$0 \to A/f^n \to A/f^{n + 1} \to A/f \to 0$ we see that
the associated primes of $A/f^nA$ agree with the associated
primes of $A/fA$. Since the associated points of $\mathcal{F}_n$
correspond to the associated primes of $A/f^nA$ not equal to $\mathfrak m$
by assumption (3), we conclude that
$M_n = H^0(U, \mathcal{F}_n)$ is a finite $A$-module by (5) and
Local Cohomology, Proposition \ref{local-cohomology-proposition-kollar}.
Thus hypothesis (c) holds.

\medskip\noindent
To finish the proof it suffices to show that there exists an $n > 1$
such that the image of
$$
H^1(U, \mathcal{F}_n) \longrightarrow H^1(U, \mathcal{F}_1)
$$
has finite length as an $A$-module. Namely, this will imply hypothesis (a)
by Cohomology, Lemma \ref{cohomology-lemma-ML-better}. The image is independent
of $n$ for $n$ large enough by Lemma \ref{lemma-ML-local}.
Let $\omega_A^\bullet$ be a normalized dualizing complex for $A$.
By the local duality theorem and Matlis duality
(Dualizing Complexes, Lemma \ref{dualizing-lemma-special-case-local-duality}
and Proposition \ref{dualizing-proposition-matlis})
our claim is equivalent to: the image of
$$
\text{Ext}^{-2}_A(M_1, \omega_A^\bullet) \to
\text{Ext}^{-2}_A(M_n, \omega_A^\bullet)
$$
has finite length for $n \gg 1$. The modules in question are
finite $A$-modules supported at $V(f)$. Thus it suffices to show that this
map is zero after localization at a prime $\mathfrak q$
containing $f$ and different from $\mathfrak m$.
Let $\omega_{A_\mathfrak q}^\bullet$ be a normalized
dualizing complex on $A_\mathfrak q$ and recall that
$\omega_{A_\mathfrak q}^\bullet =
(\omega_A^\bullet)_\mathfrak q[\dim(A/\mathfrak q)]$ by
Dualizing Complexes, Lemma \ref{dualizing-lemma-dimension-function}.
Using the local structure of $\mathcal{F}_n$ given in (4)
we find that it suffices to show the vanishing of
$$
\text{Ext}^{-2 + \dim(A/\mathfrak q)}_{A_\mathfrak q}(
A_\mathfrak q/f, \omega_{A_\mathfrak q}^\bullet)
\to
\text{Ext}^{-2 + \dim(A/\mathfrak q)}_{A_\mathfrak q}(
A_\mathfrak q/f^n, \omega_{A_\mathfrak q}^\bullet)
$$
for $n$ large enough. If $\dim(A/\mathfrak q) > 3$, then this is immediate from
Local Cohomology, Lemma \ref{local-cohomology-lemma-sitting-in-degrees}.
For the other cases we will use the long exact sequence
$$
\ldots
\xrightarrow{f^n}
H^{-1}(\omega_{A_\mathfrak q}^\bullet)
\to
\text{Ext}^{-1}_{A_\mathfrak q}(
A_\mathfrak q/f^n, \omega_{A_\mathfrak q}^\bullet) \to
H^0(\omega_{A_\mathfrak q}^\bullet)
\xrightarrow{f^n}
H^0(\omega_{A_\mathfrak q}^\bullet)
\to
\text{Ext}^0_{A_\mathfrak q}(
A_\mathfrak q/f^n, \omega_{A_\mathfrak q}^\bullet) \to 0
$$
If $\dim(A/\mathfrak q) = 2$, then
$H^0(\omega_{A_\mathfrak q}^\bullet) = 0$
because $\text{depth}(A_\mathfrak q) \geq 1$ as
$f$ is a nonzerodivisor.
Thus the long exact sequence shows the condition is that
$$
f^{n - 1} :
H^{-1}(\omega_{A_\mathfrak q}^\bullet)/f \to
H^{-1}(\omega_{A_\mathfrak q}^\bullet)/f^n
$$
is zero. Now $H^{-1}(\omega^\bullet_\mathfrak q)$ is a finite
module supported in the primes $\mathfrak p \subset A_\mathfrak q$ such that
$\text{depth}(A_\mathfrak p) + \dim((A/\mathfrak p)_\mathfrak q) \leq 1$.
Since $\dim((A/\mathfrak p)_\mathfrak q) = \dim(A/\mathfrak p) - 2$
condition (6) tells us these primes are contained in $V(f)$.
Thus the desired vanishing for $n$ large enough.
Finally, if $\dim(A/\mathfrak q) = 1$, then condition (5) combined
with the fact that $f$ is a nonzerodivisor
insures that $A_\mathfrak q$ has depth at least $2$. Hence
$H^0(\omega_{A_\mathfrak q}^\bullet) =
H^{-1}(\omega_{A_\mathfrak q}^\bullet) = 0$
and the long exact sequence shows the claim is
equivalent to the vanishing of
$$
f^{n - 1} :
H^{-2}(\omega_{A_\mathfrak q}^\bullet)/f \to
H^{-2}(\omega_{A_\mathfrak q}^\bullet)/f^n
$$
Now $H^{-2}(\omega^\bullet_\mathfrak q)$ is a finite
module supported in the primes $\mathfrak p \subset A_\mathfrak q$
such that $\text{depth}(A_\mathfrak p) + \dim((A/\mathfrak p)_\mathfrak q)
\leq 2$. By condition (6) all of these primes are contained in $V(f)$.
Thus the desired vanishing for $n$ large enough.
\end{proof}
